\chapter*{Concordance}
\addcontentsline{toc}{chapter}{Concordance}
The chapters are the papers of the series, re-ordered by subject. A reference of the form ``[Ch.~$n$, Thm.~$a.b$]'' points to Theorem $n.a.b$ of this book; the papers' own section-wise numbering is preserved inside each chapter. References to the applied notes (Phys~1, Bio~1--5, \dots) are by the labels below.

\begin{longtable}{llrp{8.2cm}}
\toprule
paper & folder & chapter & title\\
\midrule
\endhead
I & \texttt{Paper 1} & 2 & the Langlands--thermodynamic bridge over Bruhat--Tits trees\\
II & \texttt{Paper 2} & 3 & the trace obstruction and the K-homological bulk\\
VI & \texttt{Paper 6} & 4 & the shadow algebra --- a functorial Kirchberg realization of the critical group\\
III & \texttt{Paper 3} & 5 & the Iwahori--Hecke tower and the type III dissolution of the trace obstruction\\
IV & \texttt{Paper 4} & 6 & the limit module of abelian towers and the characteristic identity\\
IX & \texttt{Paper 9} & 7 & exceptional zeros, the graph $\cL$-invariant, and the operator-algebraic
Mazur--Tate--Teitelbaum formula\\
XIV & \texttt{Paper 14} & 8 & the graph $\cL$-invariant as spectral action, Bloch effective mass and Albanese length, and what the cyclic Chern character sees\\
V & \texttt{Paper 5} & 9 & the degeneracy structure of the Iwahori tower and the Euler correction\\
VII & \texttt{Paper 7} & 10 & the snake factorization of the transfer and the evaluated characteristic form\\
XIII & \texttt{Paper 13} & 11 & universal localization of the degeneracy algebra, the $K_1$-characteristic class, and the
Eisenstein boundary\\
VIII & \texttt{Paper 8} & 12 & the tower-level Hecke--Kirchberg limit and the non-commutative Iwasawa $K$-module\\
XVI & \texttt{Paper 16} & 13 & the out-covering correspondence of the head map, its $K_0$-index $\et$, and the
$\KK$-category of the Iwahori tower\\
XVII & \texttt{Paper 17} & 14 & the tail norm at composite $q$, and the splitting of the limit $K$-module\\
XIX & \texttt{Paper 19} & 15 & the odd cyclic class that does not exist, and what the exceptional zero sees instead\\
X & \texttt{Paper 10} & 16 & $\Atwo$ complexes, higher simplicial Laplacians, and the fate of the trace obstruction in
rank two\\
XI & \texttt{Paper 11} & 17 & the $\PGL_3$ positive Hecke-dynamical system: shift equivalence over $\Z$ and the geodesic
edge flow\\
XII & \texttt{Paper 12} & 18 & $\Atwo$ towers --- no abelian sector, the covering norm theorem, and the two digraph towers\\
XV & \texttt{Paper 15} & 19 & the parahoric factorization of the geodesic edge flow, the unitarity of the remainder,
and the chamber companion\\
XVIII & \texttt{Paper 18} & 20 & the unitary remainder in rank $d$\\
XX & \texttt{Paper 20} & 21 & the invariant subspace in rank $d$, and one flag count that produces it\\
Phys1 & \texttt{Phys 1} & 22 & Effective mass from spanning trees, and the absence of Bloch phases on rank-two hyperbolic lattices\\
Holo1 & \texttt{Holo 1} & 23 & What the boundary algebra of $p$-adic AdS/CFT does not see\\
Net1 & \texttt{Net 1} & 24 & The sign of $\lambda$: an algebraic certificate for how a graph tower is generated\\
Ctrl1 & \texttt{Ctrl 1} & 25 & Power-sum lower bounds on hidden eigenvalues in positive realizations\\
Bio1 & \texttt{Bio 1} & 26 & Assembly ambiguity is not determined by branch statistics: critical groups of de Bruijn graphs\\
Bio2 & \texttt{Bio 2} & 27 & The sandpile group as a coordinate system on genome reconstructions\\
Bio3 & \texttt{Bio 3} & 28 & The sandpile group of a de Bruijn graph splits along its repeats\\
Bio4 & \texttt{Bio 4} & 29 & Multi-copy repeats: where the two-copy theory of assembly ambiguity stops\\
Bio5 & \texttt{Bio 5} & 30 & The bundle-chain decomposition of a bacterial genome: \emph{Escherichia coli} K-12 at five $k$-mer lengths\\
Bio6 & \texttt{Bio 6} & 31 & Double strands without a new lemma: the reverse-complement double cover of a de Bruijn graph\\
Bio7 & \texttt{Bio 7} & 32 & The reverse-complement quotient is a signed graph: what Zaslavsky's theory gives, and what it cannot\\
\bottomrule
\end{longtable}
